\section{Mã nguồn và cách chạy lại}\label{app:code}

\paragraph{Cấu trúc.}
\begin{itemize}
\item \code{final/data.py} --- đọc dữ liệu, tạo nhãn 7 lớp và đặc trưng, chia ngẫu nhiên / theo địa chỉ, undersampling, dựng lịch sử
  địa chỉ (biểu diễn B), chuẩn hoá (chỉ tính trên train).
\item \code{final/models.py} --- \code{FeatureTokenizer} (biểu diễn A), \code{RNNNet} (LSTM/GRU, 1--2 tầng, một/hai chiều,
  last/mean/attention, \code{pack\_padded\_sequence}), \code{CNNNet} (Conv1D, BN, Dropout, Max/Avg pooling, Flatten/GAP/GMP có mặt nạ),
  \code{MLPNet}.
\item \code{final/configs.py} --- toàn bộ cấu hình thực nghiệm (tên, kiến trúc, siêu tham số).
\item \code{final/train.py} --- huấn luyện với early stopping, ghi lịch sử từng epoch, hiệu chỉnh trọng số lớp trên validation,
  đánh giá test; lưu \code{history.json}, \code{result.json}, \code{model.pt}, \code{pred\_test.npy} vào
  \code{outputs/final/<biểu diễn>\_<cách chia>/<cấu hình>/}.
\item \code{final/baselines.py} --- XGBoost (GPU) trên hai cách chia với 3 mức thông tin lịch sử.
\item \code{final/benchmark.py} --- đo tốc độ huấn luyện từng mô hình khi chạy riêng.
\item \code{final/report.py} --- tổng hợp kết quả thành các bảng LaTeX/CSV (\code{outputs/final/tables/}) và hình
  (\code{outputs/final/figures/}) dùng trong báo cáo.
\item \code{final\_project.ipynb} --- notebook trình bày đầu--cuối: thống kê dữ liệu, minh hoạ hai biểu diễn, in kiến trúc, bảng kết quả,
  đường cong huấn luyện, ma trận nhầm lẫn, dự đoán mẫu và phân tích lỗi. Mặc định notebook \emph{nạp kết quả đã huấn luyện};
  đặt \code{RETRAIN = True} để huấn luyện lại một số cấu hình tiêu biểu ngay trong notebook.
\end{itemize}

\paragraph{Cách chạy.} Cần Python 3.12, các thư viện trong \code{requirements.txt} (thêm PyTorch có CUDA nếu dùng GPU) và tệp
\code{BitcoinHeistData.csv} ở thư mục gốc. Siêu tham số XGBoost được đọc từ \code{exp\_out/best\_params.json} (kết quả tinh chỉnh
của Bài 1).
{\footnotesize
\begin{verbatim}
pip install -r requirements.txt
python -m final.train --rep A            # mọi cấu hình A, chia ngẫu nhiên
python -m final.train --rep B            # mọi cấu hình B, chia theo địa chỉ
python -m final.train --rep A --split group --only A_gru+A_lstm+A_mlp+A_cnn_bn_do
python -m final.train --rep B --seeds 43+44 --only B_gru+B_cnn_bn_do  # seed khác
python -m final.baselines                # XGBoost
python -m final.benchmark                # đo tốc độ (khi GPU rảnh)
python -m final.report                   # bảng + hình
cd report && latexmk -xelatex main.tex   # biên dịch báo cáo
\end{verbatim}
}
Có thể chạy nhiều tiến trình \code{final.train} song song với cùng tham số: mỗi tiến trình ``nhận'' một cấu hình chưa chạy
(qua một tệp khoá) nên không chạy trùng. Thời gian: $\approx$0,5--2,5 phút/cấu hình A và 0,5--7 phút/cấu hình B trên GPU RTX PRO 4000 (khi chạy song song).

\paragraph{Tái lập.} Chia dữ liệu và undersampling luôn dùng \code{random\_state=42}; seed huấn luyện (42, 43, 44) cố định khởi tạo trọng số và thứ tự batch. Do một số phép
toán cuDNN không tất định, chạy lại có thể khác ở chữ số thập phân thứ 3 --- nhỏ hơn độ dao động giữa các seed.
